% ============================================================
%  preamble.tex — Matemagica
%  Da includere in main.tex con % ============================================================
%  preamble.tex — Matemagica
%  Da includere in main.tex con % ============================================================
%  preamble.tex — Matemagica
%  Da includere in main.tex con % ============================================================
%  preamble.tex — Matemagica
%  Da includere in main.tex con \input{preamble}
% ============================================================

% --- Compilatore -----------------------------------------------
% Richiede LuaLaTeX

% --- Pacchetti di base -----------------------------------------
\usepackage{fontspec}          % Gestione font OpenType (LuaLaTeX)
\usepackage{unicode-math}      % Matematica Unicode
\usepackage{microtype}         % Ottimizzazione tipografica
\microtypesetup{verbose=silent}
\usepackage{geometry}          % Impostazione pagina
\usepackage{setspace}          % Interlinea
\usepackage{parskip}           % Spaziatura tra paragrafi
\usepackage{ragged2e}          % Testo allineato a sinistra
\usepackage{xcolor}            % Colori
\usepackage{tcolorbox}         % Riquadri colorati
\tcbuselibrary{skins, breakable}
\usepackage{booktabs}          % Tabelle professionali
\usepackage{tikz}              % Disegni e diagrammi
\usetikzlibrary{decorations.pathreplacing}
\usepackage{amsmath}           % Matematica
\usepackage{enumitem}          % Liste personalizzate
\usepackage{hyperref}          % Link PDF

% --- Font ------------------------------------------------------
\setmainfont{Atkinson Hyperlegible}
\setmathfont{Latin Modern Math}

% --- Geometria pagina (A5) -------------------------------------
\geometry{
  paper      = a4paper,
  top        = 2.0cm,
  bottom     = 2.0cm,
  inner      = 2.2cm,   % margine interno leggermente più largo
  outer      = 1.8cm,
  marginpar  = 0cm,
}

% --- Tipografia ------------------------------------------------
\setstretch{1.5}              % Interlinea aumentata (accessibilità)
\RaggedRight                  % Testo allineato a sinistra (no giustificazione)
\setlength{\parindent}{0pt}   % Nessun rientro (già gestito da parskip)
\setlength{\parskip}{6pt}

% --- Palette colori --------------------------------------------
\definecolor{blu}{HTML}{2E6DA4}        % Primario — titoli, accenti
\definecolor{azzurro}{HTML}{D6EAF8}    % Secondario — riquadri definizioni
\definecolor{ambra}{HTML}{FEF9E7}      % Terziario — riquadri attenzione
\definecolor{ambraborder}{HTML}{F0C040}% Bordo ambra
\definecolor{lightgray}{HTML}{999999}      % Testi secondari (piè di pagina)
\definecolor{darkgray}{HTML}{666666}       % Sottotitolo

% --- Stili titoli di capitolo/sezione --------------------------
\usepackage{titlesec}
\usepackage{needspace}
\titleformat{\chapter}[block]
{\huge\bfseries\color{blu}}
{\thechapter.}{0.5em}{}
\titleformat{\section}[block]
{\large\bfseries\color{blu}}
{\thesection}{0.5em}{}
\titleformat{\subsection}[block]
{\normalsize\bfseries\color{blu}}
{\thesubsection}{0.5em}{}
\titlespacing*{\section}{0pt}{1.5ex plus 0.5ex minus 0.2ex}{0.8ex}
\titlespacing*{\subsection}{0pt}{1.2ex plus 0.4ex minus 0.2ex}{0.6ex}
\preto\section{\needspace{6\baselineskip}}
\preto\subsection{\needspace{5\baselineskip}}
\preto\subsubsection{\needspace{4\baselineskip}}

% --- Intestazione e piè di pagina -----------------------------
\usepackage{datetime2}         % Data di compilazione (LuaLaTeX-safe)
\usepackage{fancyhdr}

% Formato data AAAAMMDD — definito e attivato prima di \DTMtoday
\DTMnewdatestyle{aaaammdd}{%
  \renewcommand{\DTMdisplaydate}[4]{{\small\color{gray}ver. \number##1\DTMtwodigits{##2}\DTMtwodigits{##3}}}%
  \renewcommand{\DTMDisplaydate}{\DTMdisplaydate}%
}
\DTMsetdatestyle{aaaammdd}

\pagestyle{fancy}
\fancyhf{}                     % Azzera tutti i campi

% Piè di pagina: data sul lato interno, numero pagina sul lato esterno
% Pagine dispari (destra): data a sinistra (interno), numero a destra (esterno)
% Pagine pari  (sinistra): numero a sinistra (esterno), data a destra (interno)
\fancyfoot[LO]{\DTMtoday}
\fancyfoot[RO]{\thepage}
\fancyfoot[LE]{\thepage}
\fancyfoot[RE]{\DTMtoday}

\fancyhead{}                        % Nessuna intestazione
\renewcommand{\headrulewidth}{0pt}  % Nessuna riga sotto l'intestazione
\renewcommand{\footrulewidth}{0pt}  % Nessuna riga sopra il piè di pagina

% Stile plain — applicato automaticamente alla prima pagina di ogni capitolo
\fancypagestyle{plain}{
  \fancyhf{}
  \fancyfoot[LO]{\DTMtoday}
  \fancyfoot[RO]{\thepage}
  \fancyfoot[LE]{\thepage}
  \fancyfoot[RE]{\DTMtoday}
  \renewcommand{\headrulewidth}{0pt}
}

% --- Hyperref --------------------------------------------------
\hypersetup{
  colorlinks = true,
  linkcolor  = blu,
  urlcolor   = blu,
  pdftitle   = {Matemagica},
  pdfauthor  = {},
}

% ============================================================
%  stili/riquadri.tex — Ambienti riquadro
%  (può essere separato in un file a parte)
% ============================================================

% --- Riquadro DEFINIZIONE (azzurro) ----------------------------
\newtcolorbox{definizione}{
  enhanced,
  %  breakable,
  colback    = azzurro,
  colframe   = blu,
  fonttitle  = \bfseries,
  title      = Definizione,
  left       = 6pt,
  right      = 6pt,
  top        = 4pt,
  bottom     = 4pt,
  arc        = 4pt,
}

% --- Riquadro REGOLA (azzurro, bordo più sottile) ---------------
\newtcolorbox{regola}{
  enhanced,
  %  breakable,
  colback    = azzurro,
  colframe   = blu,
  fonttitle  = \bfseries,
  title      = Regola,
  left       = 6pt,
  right      = 6pt,
  top        = 4pt,
  bottom     = 4pt,
  arc        = 4pt,
}

% --- Riquadro ESEMPIO (bianco con bordo blu) --------------------
\newtcolorbox{esempio}{
  enhanced,
  breakable,
  colback    = white,
  colframe   = blu,
  fonttitle  = \bfseries,
  title      = Esempio,
  left       = 6pt,
  right      = 6pt,
  top        = 4pt,
  bottom     = 4pt,
  arc        = 4pt,
  boxrule    = 0.6pt,
}

% --- Riquadro ATTENZIONE (ambra) --------------------------------
\newtcolorbox{attenzione}{
  enhanced,
  breakable,
  colback    = ambra,
  colframe   = ambraborder,
  left       = 6pt,
  right      = 6pt,
  top        = 4pt,
  bottom     = 4pt,
  arc        = 4pt,
}

% --- Riquadro LO SAPEVI? (blu attenuato) ------------------------
\newtcolorbox{sapevi}{
  enhanced,
  breakable,
  colback    = blu!10,
  colframe   = blu!60,
  left       = 6pt,
  right      = 6pt,
  top        = 4pt,
  bottom     = 4pt,
  arc        = 4pt,
  fontupper  = \small,
}
% ============================================================

% --- Compilatore -----------------------------------------------
% Richiede LuaLaTeX

% --- Pacchetti di base -----------------------------------------
\usepackage{fontspec}          % Gestione font OpenType (LuaLaTeX)
\usepackage{unicode-math}      % Matematica Unicode
\usepackage{microtype}         % Ottimizzazione tipografica
\microtypesetup{verbose=silent}
\usepackage{geometry}          % Impostazione pagina
\usepackage{setspace}          % Interlinea
\usepackage{parskip}           % Spaziatura tra paragrafi
\usepackage{ragged2e}          % Testo allineato a sinistra
\usepackage{xcolor}            % Colori
\usepackage{tcolorbox}         % Riquadri colorati
\tcbuselibrary{skins, breakable}
\usepackage{booktabs}          % Tabelle professionali
\usepackage{tikz}              % Disegni e diagrammi
\usetikzlibrary{decorations.pathreplacing}
\usepackage{amsmath}           % Matematica
\usepackage{enumitem}          % Liste personalizzate
\usepackage{hyperref}          % Link PDF

% --- Font ------------------------------------------------------
\setmainfont{Atkinson Hyperlegible}
\setmathfont{Latin Modern Math}

% --- Geometria pagina (A5) -------------------------------------
\geometry{
  paper      = a4paper,
  top        = 2.0cm,
  bottom     = 2.0cm,
  inner      = 2.2cm,   % margine interno leggermente più largo
  outer      = 1.8cm,
  marginpar  = 0cm,
}

% --- Tipografia ------------------------------------------------
\setstretch{1.5}              % Interlinea aumentata (accessibilità)
\RaggedRight                  % Testo allineato a sinistra (no giustificazione)
\setlength{\parindent}{0pt}   % Nessun rientro (già gestito da parskip)
\setlength{\parskip}{6pt}

% --- Palette colori --------------------------------------------
\definecolor{blu}{HTML}{2E6DA4}        % Primario — titoli, accenti
\definecolor{azzurro}{HTML}{D6EAF8}    % Secondario — riquadri definizioni
\definecolor{ambra}{HTML}{FEF9E7}      % Terziario — riquadri attenzione
\definecolor{ambraborder}{HTML}{F0C040}% Bordo ambra
\definecolor{lightgray}{HTML}{999999}      % Testi secondari (piè di pagina)
\definecolor{darkgray}{HTML}{666666}       % Sottotitolo

% --- Stili titoli di capitolo/sezione --------------------------
\usepackage{titlesec}
\usepackage{needspace}
\titleformat{\chapter}[block]
{\huge\bfseries\color{blu}}
{\thechapter.}{0.5em}{}
\titleformat{\section}[block]
{\large\bfseries\color{blu}}
{\thesection}{0.5em}{}
\titleformat{\subsection}[block]
{\normalsize\bfseries\color{blu}}
{\thesubsection}{0.5em}{}
\titlespacing*{\section}{0pt}{1.5ex plus 0.5ex minus 0.2ex}{0.8ex}
\titlespacing*{\subsection}{0pt}{1.2ex plus 0.4ex minus 0.2ex}{0.6ex}
\preto\section{\needspace{6\baselineskip}}
\preto\subsection{\needspace{5\baselineskip}}
\preto\subsubsection{\needspace{4\baselineskip}}

% --- Intestazione e piè di pagina -----------------------------
\usepackage{datetime2}         % Data di compilazione (LuaLaTeX-safe)
\usepackage{fancyhdr}

% Formato data AAAAMMDD — definito e attivato prima di \DTMtoday
\DTMnewdatestyle{aaaammdd}{%
  \renewcommand{\DTMdisplaydate}[4]{{\small\color{gray}ver. \number##1\DTMtwodigits{##2}\DTMtwodigits{##3}}}%
  \renewcommand{\DTMDisplaydate}{\DTMdisplaydate}%
}
\DTMsetdatestyle{aaaammdd}

\pagestyle{fancy}
\fancyhf{}                     % Azzera tutti i campi

% Piè di pagina: data sul lato interno, numero pagina sul lato esterno
% Pagine dispari (destra): data a sinistra (interno), numero a destra (esterno)
% Pagine pari  (sinistra): numero a sinistra (esterno), data a destra (interno)
\fancyfoot[LO]{\DTMtoday}
\fancyfoot[RO]{\thepage}
\fancyfoot[LE]{\thepage}
\fancyfoot[RE]{\DTMtoday}

\fancyhead{}                        % Nessuna intestazione
\renewcommand{\headrulewidth}{0pt}  % Nessuna riga sotto l'intestazione
\renewcommand{\footrulewidth}{0pt}  % Nessuna riga sopra il piè di pagina

% Stile plain — applicato automaticamente alla prima pagina di ogni capitolo
\fancypagestyle{plain}{
  \fancyhf{}
  \fancyfoot[LO]{\DTMtoday}
  \fancyfoot[RO]{\thepage}
  \fancyfoot[LE]{\thepage}
  \fancyfoot[RE]{\DTMtoday}
  \renewcommand{\headrulewidth}{0pt}
}

% --- Hyperref --------------------------------------------------
\hypersetup{
  colorlinks = true,
  linkcolor  = blu,
  urlcolor   = blu,
  pdftitle   = {Matemagica},
  pdfauthor  = {},
}

% ============================================================
%  stili/riquadri.tex — Ambienti riquadro
%  (può essere separato in un file a parte)
% ============================================================

% --- Riquadro DEFINIZIONE (azzurro) ----------------------------
\newtcolorbox{definizione}{
  enhanced,
  %  breakable,
  colback    = azzurro,
  colframe   = blu,
  fonttitle  = \bfseries,
  title      = Definizione,
  left       = 6pt,
  right      = 6pt,
  top        = 4pt,
  bottom     = 4pt,
  arc        = 4pt,
}

% --- Riquadro REGOLA (azzurro, bordo più sottile) ---------------
\newtcolorbox{regola}{
  enhanced,
  %  breakable,
  colback    = azzurro,
  colframe   = blu,
  fonttitle  = \bfseries,
  title      = Regola,
  left       = 6pt,
  right      = 6pt,
  top        = 4pt,
  bottom     = 4pt,
  arc        = 4pt,
}

% --- Riquadro ESEMPIO (bianco con bordo blu) --------------------
\newtcolorbox{esempio}{
  enhanced,
  breakable,
  colback    = white,
  colframe   = blu,
  fonttitle  = \bfseries,
  title      = Esempio,
  left       = 6pt,
  right      = 6pt,
  top        = 4pt,
  bottom     = 4pt,
  arc        = 4pt,
  boxrule    = 0.6pt,
}

% --- Riquadro ATTENZIONE (ambra) --------------------------------
\newtcolorbox{attenzione}{
  enhanced,
  breakable,
  colback    = ambra,
  colframe   = ambraborder,
  left       = 6pt,
  right      = 6pt,
  top        = 4pt,
  bottom     = 4pt,
  arc        = 4pt,
}

% --- Riquadro LO SAPEVI? (blu attenuato) ------------------------
\newtcolorbox{sapevi}{
  enhanced,
  breakable,
  colback    = blu!10,
  colframe   = blu!60,
  left       = 6pt,
  right      = 6pt,
  top        = 4pt,
  bottom     = 4pt,
  arc        = 4pt,
  fontupper  = \small,
}
% ============================================================

% --- Compilatore -----------------------------------------------
% Richiede LuaLaTeX

% --- Pacchetti di base -----------------------------------------
\usepackage{fontspec}          % Gestione font OpenType (LuaLaTeX)
\usepackage{unicode-math}      % Matematica Unicode
\usepackage{microtype}         % Ottimizzazione tipografica
\microtypesetup{verbose=silent}
\usepackage{geometry}          % Impostazione pagina
\usepackage{setspace}          % Interlinea
\usepackage{parskip}           % Spaziatura tra paragrafi
\usepackage{ragged2e}          % Testo allineato a sinistra
\usepackage{xcolor}            % Colori
\usepackage{tcolorbox}         % Riquadri colorati
\tcbuselibrary{skins, breakable}
\usepackage{booktabs}          % Tabelle professionali
\usepackage{tikz}              % Disegni e diagrammi
\usetikzlibrary{decorations.pathreplacing}
\usepackage{amsmath}           % Matematica
\usepackage{enumitem}          % Liste personalizzate
\usepackage{hyperref}          % Link PDF

% --- Font ------------------------------------------------------
\setmainfont{Atkinson Hyperlegible}
\setmathfont{Latin Modern Math}

% --- Geometria pagina (A5) -------------------------------------
\geometry{
  paper      = a4paper,
  top        = 2.0cm,
  bottom     = 2.0cm,
  inner      = 2.2cm,   % margine interno leggermente più largo
  outer      = 1.8cm,
  marginpar  = 0cm,
}

% --- Tipografia ------------------------------------------------
\setstretch{1.5}              % Interlinea aumentata (accessibilità)
\RaggedRight                  % Testo allineato a sinistra (no giustificazione)
\setlength{\parindent}{0pt}   % Nessun rientro (già gestito da parskip)
\setlength{\parskip}{6pt}

% --- Palette colori --------------------------------------------
\definecolor{blu}{HTML}{2E6DA4}        % Primario — titoli, accenti
\definecolor{azzurro}{HTML}{D6EAF8}    % Secondario — riquadri definizioni
\definecolor{ambra}{HTML}{FEF9E7}      % Terziario — riquadri attenzione
\definecolor{ambraborder}{HTML}{F0C040}% Bordo ambra
\definecolor{lightgray}{HTML}{999999}      % Testi secondari (piè di pagina)
\definecolor{darkgray}{HTML}{666666}       % Sottotitolo

% --- Stili titoli di capitolo/sezione --------------------------
\usepackage{titlesec}
\usepackage{needspace}
\titleformat{\chapter}[block]
{\huge\bfseries\color{blu}}
{\thechapter.}{0.5em}{}
\titleformat{\section}[block]
{\large\bfseries\color{blu}}
{\thesection}{0.5em}{}
\titleformat{\subsection}[block]
{\normalsize\bfseries\color{blu}}
{\thesubsection}{0.5em}{}
\titlespacing*{\section}{0pt}{1.5ex plus 0.5ex minus 0.2ex}{0.8ex}
\titlespacing*{\subsection}{0pt}{1.2ex plus 0.4ex minus 0.2ex}{0.6ex}
\preto\section{\needspace{6\baselineskip}}
\preto\subsection{\needspace{5\baselineskip}}
\preto\subsubsection{\needspace{4\baselineskip}}

% --- Intestazione e piè di pagina -----------------------------
\usepackage{datetime2}         % Data di compilazione (LuaLaTeX-safe)
\usepackage{fancyhdr}

% Formato data AAAAMMDD — definito e attivato prima di \DTMtoday
\DTMnewdatestyle{aaaammdd}{%
  \renewcommand{\DTMdisplaydate}[4]{{\small\color{gray}ver. \number##1\DTMtwodigits{##2}\DTMtwodigits{##3}}}%
  \renewcommand{\DTMDisplaydate}{\DTMdisplaydate}%
}
\DTMsetdatestyle{aaaammdd}

\pagestyle{fancy}
\fancyhf{}                     % Azzera tutti i campi

% Piè di pagina: data sul lato interno, numero pagina sul lato esterno
% Pagine dispari (destra): data a sinistra (interno), numero a destra (esterno)
% Pagine pari  (sinistra): numero a sinistra (esterno), data a destra (interno)
\fancyfoot[LO]{\DTMtoday}
\fancyfoot[RO]{\thepage}
\fancyfoot[LE]{\thepage}
\fancyfoot[RE]{\DTMtoday}

\fancyhead{}                        % Nessuna intestazione
\renewcommand{\headrulewidth}{0pt}  % Nessuna riga sotto l'intestazione
\renewcommand{\footrulewidth}{0pt}  % Nessuna riga sopra il piè di pagina

% Stile plain — applicato automaticamente alla prima pagina di ogni capitolo
\fancypagestyle{plain}{
  \fancyhf{}
  \fancyfoot[LO]{\DTMtoday}
  \fancyfoot[RO]{\thepage}
  \fancyfoot[LE]{\thepage}
  \fancyfoot[RE]{\DTMtoday}
  \renewcommand{\headrulewidth}{0pt}
}

% --- Hyperref --------------------------------------------------
\hypersetup{
  colorlinks = true,
  linkcolor  = blu,
  urlcolor   = blu,
  pdftitle   = {Matemagica},
  pdfauthor  = {},
}

% ============================================================
%  stili/riquadri.tex — Ambienti riquadro
%  (può essere separato in un file a parte)
% ============================================================

% --- Riquadro DEFINIZIONE (azzurro) ----------------------------
\newtcolorbox{definizione}{
  enhanced,
  %  breakable,
  colback    = azzurro,
  colframe   = blu,
  fonttitle  = \bfseries,
  title      = Definizione,
  left       = 6pt,
  right      = 6pt,
  top        = 4pt,
  bottom     = 4pt,
  arc        = 4pt,
}

% --- Riquadro REGOLA (azzurro, bordo più sottile) ---------------
\newtcolorbox{regola}{
  enhanced,
  %  breakable,
  colback    = azzurro,
  colframe   = blu,
  fonttitle  = \bfseries,
  title      = Regola,
  left       = 6pt,
  right      = 6pt,
  top        = 4pt,
  bottom     = 4pt,
  arc        = 4pt,
}

% --- Riquadro ESEMPIO (bianco con bordo blu) --------------------
\newtcolorbox{esempio}{
  enhanced,
  breakable,
  colback    = white,
  colframe   = blu,
  fonttitle  = \bfseries,
  title      = Esempio,
  left       = 6pt,
  right      = 6pt,
  top        = 4pt,
  bottom     = 4pt,
  arc        = 4pt,
  boxrule    = 0.6pt,
}

% --- Riquadro ATTENZIONE (ambra) --------------------------------
\newtcolorbox{attenzione}{
  enhanced,
  breakable,
  colback    = ambra,
  colframe   = ambraborder,
  left       = 6pt,
  right      = 6pt,
  top        = 4pt,
  bottom     = 4pt,
  arc        = 4pt,
}

% --- Riquadro LO SAPEVI? (blu attenuato) ------------------------
\newtcolorbox{sapevi}{
  enhanced,
  breakable,
  colback    = blu!10,
  colframe   = blu!60,
  left       = 6pt,
  right      = 6pt,
  top        = 4pt,
  bottom     = 4pt,
  arc        = 4pt,
  fontupper  = \small,
}
% ============================================================

% --- Compilatore -----------------------------------------------
% Richiede LuaLaTeX

% --- Pacchetti di base -----------------------------------------
\usepackage{fontspec}          % Gestione font OpenType (LuaLaTeX)
\usepackage{unicode-math}      % Matematica Unicode
\usepackage{microtype}         % Ottimizzazione tipografica
\microtypesetup{verbose=silent}
\usepackage{geometry}          % Impostazione pagina
\usepackage{setspace}          % Interlinea
\usepackage{parskip}           % Spaziatura tra paragrafi
\usepackage{ragged2e}          % Testo allineato a sinistra
\usepackage{xcolor}            % Colori
\usepackage{tcolorbox}         % Riquadri colorati
\tcbuselibrary{skins, breakable}
\usepackage{booktabs}          % Tabelle professionali
\usepackage{tikz}              % Disegni e diagrammi
\usetikzlibrary{decorations.pathreplacing}
\usepackage{amsmath}           % Matematica
\usepackage{enumitem}          % Liste personalizzate
\usepackage{hyperref}          % Link PDF

% --- Font ------------------------------------------------------
\setmainfont{Atkinson Hyperlegible}
\setmathfont{Latin Modern Math}

% --- Geometria pagina (A5) -------------------------------------
\geometry{
  paper      = a4paper,
  top        = 2.0cm,
  bottom     = 2.0cm,
  inner      = 2.2cm,   % margine interno leggermente più largo
  outer      = 1.8cm,
  marginpar  = 0cm,
}

% --- Tipografia ------------------------------------------------
\setstretch{1.5}              % Interlinea aumentata (accessibilità)
\RaggedRight                  % Testo allineato a sinistra (no giustificazione)
\setlength{\parindent}{0pt}   % Nessun rientro (già gestito da parskip)
\setlength{\parskip}{6pt}

% --- Palette colori --------------------------------------------
\definecolor{blu}{HTML}{2E6DA4}        % Primario — titoli, accenti
\definecolor{azzurro}{HTML}{D6EAF8}    % Secondario — riquadri definizioni
\definecolor{ambra}{HTML}{FEF9E7}      % Terziario — riquadri attenzione
\definecolor{ambraborder}{HTML}{F0C040}% Bordo ambra
\definecolor{lightgray}{HTML}{999999}      % Testi secondari (piè di pagina)
\definecolor{darkgray}{HTML}{666666}       % Sottotitolo

% --- Stili titoli di capitolo/sezione --------------------------
\usepackage{titlesec}
\usepackage{needspace}
\titleformat{\chapter}[block]
{\huge\bfseries\color{blu}}
{\thechapter.}{0.5em}{}
\titleformat{\section}[block]
{\large\bfseries\color{blu}}
{\thesection}{0.5em}{}
\titleformat{\subsection}[block]
{\normalsize\bfseries\color{blu}}
{\thesubsection}{0.5em}{}
\titlespacing*{\section}{0pt}{1.5ex plus 0.5ex minus 0.2ex}{0.8ex}
\titlespacing*{\subsection}{0pt}{1.2ex plus 0.4ex minus 0.2ex}{0.6ex}
\preto\section{\needspace{6\baselineskip}}
\preto\subsection{\needspace{5\baselineskip}}
\preto\subsubsection{\needspace{4\baselineskip}}

% --- Intestazione e piè di pagina -----------------------------
\usepackage{datetime2}         % Data di compilazione (LuaLaTeX-safe)
\usepackage{fancyhdr}

% Formato data AAAAMMDD — definito e attivato prima di \DTMtoday
\DTMnewdatestyle{aaaammdd}{%
  \renewcommand{\DTMdisplaydate}[4]{{\small\color{gray}ver. \number##1\DTMtwodigits{##2}\DTMtwodigits{##3}}}%
  \renewcommand{\DTMDisplaydate}{\DTMdisplaydate}%
}
\DTMsetdatestyle{aaaammdd}

\pagestyle{fancy}
\fancyhf{}                     % Azzera tutti i campi

% Piè di pagina: data sul lato interno, numero pagina sul lato esterno
% Pagine dispari (destra): data a sinistra (interno), numero a destra (esterno)
% Pagine pari  (sinistra): numero a sinistra (esterno), data a destra (interno)
\fancyfoot[LO]{\DTMtoday}
\fancyfoot[RO]{\thepage}
\fancyfoot[LE]{\thepage}
\fancyfoot[RE]{\DTMtoday}

\fancyhead{}                        % Nessuna intestazione
\renewcommand{\headrulewidth}{0pt}  % Nessuna riga sotto l'intestazione
\renewcommand{\footrulewidth}{0pt}  % Nessuna riga sopra il piè di pagina

% Stile plain — applicato automaticamente alla prima pagina di ogni capitolo
\fancypagestyle{plain}{
  \fancyhf{}
  \fancyfoot[LO]{\DTMtoday}
  \fancyfoot[RO]{\thepage}
  \fancyfoot[LE]{\thepage}
  \fancyfoot[RE]{\DTMtoday}
  \renewcommand{\headrulewidth}{0pt}
}

% --- Hyperref --------------------------------------------------
\hypersetup{
  colorlinks = true,
  linkcolor  = blu,
  urlcolor   = blu,
  pdftitle   = {Matemagica},
  pdfauthor  = {},
}

% ============================================================
%  stili/riquadri.tex — Ambienti riquadro
%  (può essere separato in un file a parte)
% ============================================================

% --- Riquadro DEFINIZIONE (azzurro) ----------------------------
\newtcolorbox{definizione}{
  enhanced,
  %  breakable,
  colback    = azzurro,
  colframe   = blu,
  fonttitle  = \bfseries,
  title      = Definizione,
  left       = 6pt,
  right      = 6pt,
  top        = 4pt,
  bottom     = 4pt,
  arc        = 4pt,
}

% --- Riquadro REGOLA (azzurro, bordo più sottile) ---------------
\newtcolorbox{regola}{
  enhanced,
  %  breakable,
  colback    = azzurro,
  colframe   = blu,
  fonttitle  = \bfseries,
  title      = Regola,
  left       = 6pt,
  right      = 6pt,
  top        = 4pt,
  bottom     = 4pt,
  arc        = 4pt,
}

% --- Riquadro ESEMPIO (bianco con bordo blu) --------------------
\newtcolorbox{esempio}{
  enhanced,
  breakable,
  colback    = white,
  colframe   = blu,
  fonttitle  = \bfseries,
  title      = Esempio,
  left       = 6pt,
  right      = 6pt,
  top        = 4pt,
  bottom     = 4pt,
  arc        = 4pt,
  boxrule    = 0.6pt,
}

% --- Riquadro ATTENZIONE (ambra) --------------------------------
\newtcolorbox{attenzione}{
  enhanced,
  breakable,
  colback    = ambra,
  colframe   = ambraborder,
  left       = 6pt,
  right      = 6pt,
  top        = 4pt,
  bottom     = 4pt,
  arc        = 4pt,
}

% --- Riquadro LO SAPEVI? (blu attenuato) ------------------------
\newtcolorbox{sapevi}{
  enhanced,
  breakable,
  colback    = blu!10,
  colframe   = blu!60,
  left       = 6pt,
  right      = 6pt,
  top        = 4pt,
  bottom     = 4pt,
  arc        = 4pt,
  fontupper  = \small,
}